\addchap{Preface}

This book began with a list of textbooks that should exist but do not. One entry read, in effect: physics students all learn the harmonic oscillator, the hydrogen atom and the Ising model, and they also learn where each one fails; deep learning theory has accumulated a comparable set of solvable models, but they are scattered across dozens of papers whose notations do not talk to each other, and nobody has written down, in one place, what each of them gets wrong.

The book that resulted has sixteen chapters, each built around one phenomenon of trained networks and one model that reproduces it. Each chapter states the phenomenon with its evidence, defines the model, proves its main theorem completely, compares the solution with simulations, and then says what the model gets wrong: what it threw away, where its predictions fail, and which successor models repair it. The phenomena range from double descent and implicit bias to grokking, superposition, in-context learning, the edge of stability and scaling laws. All chapters share one notation, summarized on the next pages, and each chapter ends by translating the notation of its main sources into that of the book.

The intended readers are graduate students and researchers in machine learning, statistics, applied mathematics, theoretical computer science and physics, and practitioners who want to know which parts of deep learning are understood, in what sense, and how far the understanding reaches. The prerequisites are undergraduate analysis, linear algebra and probability. Each chapter builds what it uses, and an appendix collects the tools that several chapters share. The prologue explains the plan of the book and solves its first model.

A solvable model is a lens, not a mirror. The models in this book are chosen because they are small enough to be solved, and a reader should come away knowing both what each lens brings into focus and what it leaves blurred. If the book succeeds, the second kind of knowledge will be as precise as the first.

\paragraph{A note on authorship.} This book was written by Claude (Opus 5.5), an AI model made by Anthropic, and conceived and edited by Xiaoyu Li, who proposed the subject and its plan. Every chapter's proofs were re-derived by at least one independent referee pass, its numerical claims were checked against simulations, and its references were checked against sources wherever these could be reached; the colophon at the end of the book describes how it was made. Mistakes surely remain. Readers who find one are invited to treat it as the book treats its models: as a precise statement of what is wrong, and a pointer to what comes next.
